\section*{6. 关键技术细节}
\subsection*{6.1 下界}
\(\Omega(m)\) 下界来自输入访问。若算法没有检查某条可能影响距离顺序的边，就可以通过调整这条边的长度制造一个与算法输出冲突的实例。

\(\Omega(\log D)\) 是决策树下界。固定图可能产生 \(D\) 个 distance order，对每个 order 都能构造边权使其成为目标输出。正确算法必须区分这些叶子，因此比较深度至少为 \(\log D\)。

\(\Omega(F-n+1)\) 的作用是补上 \(\log D\) 之外的前向弧信息。这里不能简单理解成每条额外前向弧都必然对应一次比较；更准确地说，\(F-n+1\) 是论文在比较下界和 Dijkstra 比较次数分析中共同使用的图参数，用来度量基本可达性之外的前向弧自由度。

这个下界的核心是变量数和约束数的比较。选取一个有 \(F\) 条 forward arcs 的 distance order \(L=[v_1,\ldots,v_n]\)，并给前向弧 \(v_i v_j\) 赋长度 \(j-i\)。若某个算法只做至多 \(F-n\) 次比较，那么这些比较最多固定 \(F-n\) 个线性约束；再加入 \(n-1\) 个相邻顶点距离差约束，约束总数仍少于 \(F\) 个前向弧长度变量。于是还存在一个非零扰动方向，可以保持算法见到的比较结果不变，同时逐步改变某些前向弧长度。沿这个方向移动到某条非相邻前向弧长度变为零时，算法的比较记录没有变化，但原顺序 \(L\) 已不再是正确的 true distance order；算法若仍输出 \(L\)，就会出错。

\subsection*{6.2 Dijkstra 与 working-set bound}
Dijkstra 的扫描顺序确实是 true distance order。证明可以用通常的反证法：设当前被 \texttt{delete-\allowbreak{}min} 取出的顶点是 \(u\)，但仍有未扫描顶点 \(v\) 满足 \(d(v)<d(u)\)。取一条从 \(s\) 到 \(v\) 的最短路，并令 \(x\) 是这条路上第一个尚未扫描的顶点，\(y\) 是它的前驱。由于 \(y\) 已经扫描过，松弛边 \((y,x)\) 后有 \(\hat d(x)\le d(x)\le d(v)<d(u)\le \hat d(u)\)，这与 \(u\) 是队列中估计距离最小的顶点矛盾。

运行中，堆里保存 labeled 但尚未 scanned 的顶点。在固定的扫描顺序下，只有从较早扫描顶点指向较晚扫描顶点的弧，才可能参与对后续标签的有效更新。论文把每个非源点所需的一条基本前驱弧与其余前向弧分开，后者进入 \(F-n+1\) 项。这里不能理解成每条额外前向弧都机械对应一次比较；更准确地说，\(F-n+1\) 是下界和 Dijkstra 比较次数分析中共同使用的图参数。

\texttt{delete-\allowbreak{}min} 的分析更细。设 Dijkstra 的插入序列为 \(v_1,\ldots,v_n\)。对 \(v_i\)，令 \(a_i=i\)，令 \(b_i\) 表示它被删除时已经插入的最后一个顶点下标，那么 \(v_i\) 的堆生命周期就是区间 \([a_i,b_i]\)，并且 \(W(v_i)=b_i-a_i+1\)。\texttt{decrease-\allowbreak{}key} 只改变 key，不改变这个区间的左右端点。

这些区间构成一个 interval DAG \(I\)：若区间 \(i\) 在区间 \(j\) 开始前已经结束，即 \(b_i<a_j\)，就加弧 \(i\to j\)。论文引用的区间 DAG 引理给出 \(\sum_i \log(b_i-a_i+1)=\mathrm{O}(\log \#\mathrm{Top}(I))\)， 其中 \(\#\mathrm{Top}(I)\) 是 \(I\) 的拓扑序数量。剩下要做的是把 \(\#\mathrm{Top}(I)\) 接到图本身的 \(D\)。

令 \(T\) 为 Dijkstra 第一次标记每个顶点时形成的树。如果树边是 \(v_i\to v_j\)，那么 \(v_i\) 已经被删除后，\(v_j\) 才因扫描这条边而第一次插入堆，所以 \(b_i<a_j\)，这条树边也是 interval DAG 中的边。于是，\(I\) 的任意拓扑序都尊重 \(T\)。根据论文第 3 节的判定，一个有根树的任意拓扑序都是原图的 distance order，因此 \(\#\mathrm{Top}(I)\le D\)。合在一起就得到 \(\sum_v \log W(v)=\mathrm{O}(\log D)\)。这一步把 working-set heap 的局部删除代价转化成了固定图可产生的顺序数量。

\subsection*{6.3 Lookahead}
Lookahead 版本先按无权层数处理图。这里 \(\ell(v)\) 是从 \(s\) 到 \(v\) 的路径中最少顶点数；若某一层只有一个顶点，它就是 bottleneck，并且会支配所有更高层顶点。Marked bottleneck 指下一层至少有两个顶点的 bottleneck；unmarked bottleneck 则是最高层顶点，或下一层仍然只有一个顶点的 bottleneck。

算法只把非 bottleneck 放进堆 \(H\)。未处理的 bottleneck 按层数存入数组 \(B\)，\(B\) 一直取到第一个 marked bottleneck 为止。扫描普通顶点时，边松弛仍按 Dijkstra 的方式更新标签；若被更新的是非 bottleneck，就插入 \(H\) 或执行 \texttt{decrease-\allowbreak{}key}，若是 bottleneck，则只更新它在 \(B\) 中的距离记录。

主循环比较两类候选：\(H\) 中最小非 bottleneck 的 tentative distance，以及 \(B\) 中当前 bottleneck 段能够推出的距离。若堆顶更小，就删除并扫描这个非 bottleneck。否则，算法沿 \(B\) 从低层到高层扫描；对连续 unmarked bottleneck，论文用 Lemma 8.3 说明下一个 bottleneck 的真实距离可由前一个 bottleneck 和相应出边得到。若整段 \(B\) 的最大距离仍不超过堆顶，就把整段加入输出序列；否则只加入 \(B\) 中距离不超过堆顶的最长前缀。

正确性可以看成一个安全性引理：当算法把某个 bottleneck \(b\) 加入输出序列时，所有尚未输出的非 bottleneck 顶点 \(x\) 都满足 \(d(b)\le d(x)\)。原因是低层 bottleneck 支配更高层区域，堆顶又代表所有非 bottleneck 候选中的最小 tentative distance；算法只有在显式比较后才推进 bottleneck 段。因此提前输出这些 bottleneck 不会打乱 true distance order。

复杂度还要解释为什么堆里剩下的非 bottleneck 不会太多。设图中有 \(b\) 个 bottleneck，level 总数为 \(r\)。每个不含 bottleneck 的 level 至少有两个顶点，所以 \(r\le (n+b)/2\)。取一棵 BFS 树后，在同一 level 内任意排列顶点，可以得到至少 \(\prod_i |V_i|!\ge 2^{n-r}\) 个树的拓扑序，也就是原图的 distance orders。于是 \(\log D\ge n-r\ge (n-b)/2\)，非 bottleneck 顶点数 \(n-b\) 本身就是 \(\mathrm{O}(\log D)\)。这一步支付了普通堆中非 bottleneck 的基本 \texttt{delete-\allowbreak{}min} 成本。

还有两项记账容易被略过。对 \texttt{delete-\allowbreak{}min}，论文把 bottleneck 也虚拟地看成“立即插入、立即删除”的元素，再复用 working-set bound 的区间论证，得到总成本仍为 \(\mathrm{O}(\log D)\)。对数组 \(B\) 中的指数搜索和二分搜索，给每个非 bottleneck 顶点 \(v\) 记下它跨过的 bottleneck 集合 \(B(v)\)；不同选择会产生至少 \(\prod_v (|B(v)|+1)\) 个 distance orders，因此搜索成本 \(\sum_v \log(|B(v)|+1)\) 也能由 \(\mathrm{O}(\log D)\) 控制。

\subsection*{6.4 Recursive Dijkstra}
Recursive Dijkstra 把 bottleneck 当作递归边界。初始化时，所有 bottleneck 已按层数放入列表 \(L\)。调用 \(\operatorname{Dijkstra}(u)\) 时，算法为这个源点建立一个局部堆，先扫描 \(u\)，然后反复删除局部堆的最小顶点；如果删出的顶点是 bottleneck，就以它为新源点递归调用 \(\operatorname{Dijkstra}(v)\)，否则扫描它并把它插入最终列表 \(L\)。

这里的递归不是在整张图上反复重跑 Dijkstra。算法维护全局的 scanned/\allowbreak{}labeled 状态，每个顶点只扫描一次，相关出边也只由负责它的那次扫描处理。多个局部堆表示的是被 bottleneck 切开的若干前沿；较高 level 的 bottleneck 所对应的递归运行在扫描边时，可能改进一个已经存放在较低 level 递归堆中的顶点。此时算法直接在该顶点当前所属的堆中执行 \texttt{decrease-\allowbreak{}key}，而不是把它移动到当前堆。这样递归深度不会把 \(m\) 乘起来。

维护 \(L\) 要用 finger search tree。正确性上，非 bottleneck 顶点 \(v\) 被扫描时，它的父亲 \(p(v)\) 已在 \(L\) 中，且 \(d(p(v))\le d(v)\)。算法从 \(p(v)\) 附近向后搜索第一个距离至少为 \(d(v)\) 的顶点，把 \(v\) 插在它前面；前后距离关系保证了 \(L\) 仍按真实距离非递减排列。

复杂度上，令 \(v_1,\ldots,v_n\) 为顶点的扫描顺序。若 \(p(v_i)=v_j\)，则为 \(v_i\) 定义区间 \([j+1,i]\)。Finger search 的秩距离至多由这个区间长度控制，所以插入代价是 \(\mathrm{O}(1+\log(i-j))\)。再对这些区间应用 Lemma 7.1，就能把所有插入成本控制在 \(\mathrm{O}(\log D)\)。

多个局部堆的 \texttt{delete-\allowbreak{}min} 成本也要合并计算。论文主要考虑 marked bottleneck 对应的递归调用；对每个这样的 bottleneck \(v\)，取搜索树中由该调用负责的局部子树 \(T(v)\)。该局部运行的 working-set 成本可由 \(\mathrm{O}(\log \#\mathrm{Top}(T(v)))\) 控制。全局搜索树的拓扑序数量至少包含这些局部子树拓扑序数量的乘积，取对数后，各局部成本可以相加，最终仍被 \(\mathrm{O}(\log D)\) 吸收。

\subsection*{6.5 Working-set heap}
Outer heap 由一串 inner heaps 组成：

\(H_1,H_2,\ldots,H_k\)。

编号小的 heap 放较新元素，编号大的 heap 放较旧元素。每个 inner heap 可用 Fibonacci heap、hollow heap 或 rank-pairing heap 这类 fast heap 实现。插入新元素时先建立 \(H_0\)，再检查相邻层是否可以 meld。大小阈值采用双指数形式，大致保持 \(|H_i|\le 2^{2^i}\)，从而保证层数只有 \(\mathrm{O}(\log\log n)\)。

Outer heap 对外不提供 \texttt{meld}；下表中的 meld 只指 \texttt{insert} 过程中对两个 inner fast heaps 的内部合并。

各操作的运转方式可以概括如下：

\begin{center}
\footnotesize
\renewcommand{\arraystretch}{1.28}
\begin{tabularx}{\textwidth}{>{\hsize=.95\hsize}Y >{\hsize=1.45\hsize}Y >{\hsize=.60\hsize}Y}
\toprule
\textbf{操作} & \textbf{做法} & \textbf{摊还代价} \\
\midrule
\texttt{insert(x)} & 新建最年轻层，必要时按大小阈值逐层 meld 和重新编号 & \(\mathrm{O}(1)\) \\
\texttt{decrease-\allowbreak{}key(x)} & 用 union-find 定位 \(x\) 所在 inner heap，在内部减键并更新外层最小值信息 & \(\mathrm{O}(1)\) \\
\texttt{find-\allowbreak{}min()} & 通过 suffix-minimum bit vector 找到含全局最小 key 的 inner heap & \(\mathrm{O}(1)\) \\
\texttt{delete-\allowbreak{}min(x)} & 在对应 inner heap 中删除 \(x\)，再更新该层及更年轻层的 suffix minimum 标记 & \(\mathrm{O}(\log W(x))\) \\
内部 \texttt{meld} & 仅在 \texttt{insert} 中合并满足大小条件的相邻 inner heaps，并同步执行一次集合 \texttt{unite} & 计入 \texttt{insert} 的摊还代价 \\
\bottomrule
\end{tabularx}
\end{center}
\texttt{delete-\allowbreak{}min} 不会把 inner heap 分裂，只是在内部删除元素并修正外层索引；真正改变分层结构的主要是后续插入触发的内部 meld 和 reindex。Inner heaps 合并时，对应的元素集合执行一次 \texttt{unite}；\texttt{decrease-\allowbreak{}key(x)} 通过 \texttt{find(x)} 找到 \(x\) 所属的 inner heap。论文不在元素删除时同步拆分集合。对 Dijkstra 产生的操作序列，查找过程中产生的额外父指针变化可摊入该元素最终的 \(\mathrm{O}(\log W(x))\) 删除代价；在允许元素长期不被删除的一般情形下，作者再使用 fixed-tree set union 获得严格的 \(\mathrm{O}(1)\) 摊还界。Suffix-minimum bit vector 的常数时间依赖 word-RAM 操作，以及 inner heap 数量只有 \(\mathrm{O}(\log\log n)\) 这一事实；若 \(n\) 事先未知，论文还用重建技巧维持空间和字长假设。